---
format: yohaku/0.2
target: text
macro: '[[macros/raw]]'
---

## 正文
````text
#!/usr/bin/env python3
"""Fixed development actions, invoked only inside the mandatory tool sandbox."""
import os
from pathlib import Path
import signal
import shutil
import subprocess
import sys
import time

PROJECT = Path('/world/files/yohaku-default')
CARGO = '/opt/yohaku/rust/bin/cargo'
TARGET = Path('/world/tools/state/yohaku-development/target')


def commands(action):
    prefix = [CARGO]
    if action == 'check': return [prefix+['check','--offline','--locked','--all-targets']]
    if action == 'test':
        return [prefix+['test','--offline','--locked','--lib','--test','workspace','--','--test-threads=1'],
                prefix+['test','--offline','--locked','--example','yohaku_compositor','tool_windows::tests','--','--test-threads=1'],
                prefix+['test','--offline','--locked','--example','yohaku_compositor','agent_panel::entry_tests::entry_','--','--test-threads=1'],
                prefix+['test','--offline','--locked','--example','yohaku_compositor','window_protocol_tests','--','--test-threads=1']]
    if action == 'build': return [prefix+['build','--offline','--locked','--example','yohaku_compositor','--example','yohaku_terminal_window','--bin','yohaku_workspace']]
    raise ValueError('Unknown development action')


def run_command(argv, cwd, timeout=900):
    print('Yohaku development command: '+repr(argv),flush=True)
    child = subprocess.Popen(argv,cwd=cwd,start_new_session=True)
    def stop(*_):
        if child.poll() is None:
            try: os.killpg(child.pid,signal.SIGKILL)
            except ProcessLookupError: pass
    old = {n:signal.signal(n,stop) for n in (signal.SIGTERM,signal.SIGINT)}
    try:
        try: return child.wait(timeout=timeout)
        except subprocess.TimeoutExpired:
            print('Development action exceeded its wall-clock limit',flush=True);stop();child.wait();return 124
    finally:
        stop()
        for n,handler in old.items():signal.signal(n,handler)


def main():
    if os.environ.get('YOHAKU_DEVELOPMENT') != '1' or len(sys.argv)!=2:
        print('Start this action through Yohaku Techne',flush=True);return 2
    action=sys.argv[1]
    if not (PROJECT/'Cargo.toml').is_file():
        print('Development source is missing; click Develop in Techne',flush=True);return 2
    if action == 'preview':
        candidate=TARGET/'debug/examples/yohaku_compositor'
        if not candidate.is_file(): print('Build the candidate before preview',flush=True);return 2
        env=dict(os.environ,YOHAKU_WORKSPACE=str(PROJECT/'worlds/default'),YOHAKU_LOOM=str(PROJECT/'bin/loom-linux-x86_64'))
        # Candidate code receives only the existing private Wayland display and world.
        os.chdir(PROJECT);os.execve(candidate,[str(candidate)],env)
    try:
        for argv in commands(action):
            code=run_command(argv,PROJECT)
            if code: print('Development action failed: exit '+str(code),flush=True);return code
    except (OSError,ValueError) as error:
        print('Development action failed: '+str(error),flush=True);return 2
    if action=='build':
        # Read-only source/binary IME gate; no host display or bus in sandbox.
        code=run_command(['/usr/bin/python3','-B',str(PROJECT/'agent/yohaku_ime_check.py'),'--project',str(PROJECT),'--binary',str(TARGET/'debug/examples/yohaku_compositor')],PROJECT)
        if code:
            print('Candidate IME deployment gate failed: exit '+str(code),flush=True);return code
        # The candidate's Python/CLI bridge must use its newly built Rust dispatcher.
        try:shutil.copy2(TARGET/'debug/yohaku_workspace', PROJECT/'bin/yohaku_workspace')
        except OSError as error:print('Candidate bridge update failed: '+str(error),flush=True);return 2
        print('Candidate: '+str(TARGET/'debug/examples/yohaku_compositor'),flush=True)
    print('Development action passed: '+action,flush=True)
    return 0

if __name__=='__main__':sys.exit(main())
````

## 输出与关联
```yaml
output: "research/workflow-runtime/agent/yohaku_dev_runtime.py"
```
